\section{Introduction}\label{sec:id-intro}
\owner{Group}
\pages{2}

The COS212 app is a prototype tutor for the binary search tree (BST). It is designed for COS212 students to use before the lecture in which trees are taught. Its job is to let a student meet the ordering rule, practise search and insertion by hand, and receive support when they struggle, all without an instructor and without any prior knowledge of trees.

This document explains the instructional choices behind the app. It sets out the learning objectives, the learning theories applied, how the content and assessment are designed, how the app keeps students motivated and tracks their progress, and how the app will be evaluated. It also reflects on the choices the group made and records the use of generative AI, as the examination specification requires. The software engineering decisions that make these choices possible are documented separately in the companion SE document.

Two features of the target learners shape every choice in this document. The first is that they are meeting trees for the first time, so the app cannot assume any vocabulary or any sense of why the ordering rule matters. The second is that they are studying before a lecture, so the app must work as a self-contained learning experience, with no instructor available to explain a mistake. The design therefore puts explanation, feedback and practice inside the app itself.

The scope matches the SE document. The app teaches the ordering rule, the search and insertion operations, and in-order traversal. Deletion and self-balancing trees are not taught, and the reasons are given in Section~\ref{sec:id-reflection}.
